
\documentclass{article}

\usepackage{a4wide}
\usepackage[utf8]{inputenc}
\usepackage[russian]{babel}
\usepackage{graphicx}
\usepackage{amsmath}
\usepackage{float}
\usepackage{hyperref}
\usepackage{amssymb}
\usepackage{amsthm}
\usepackage{amsmath}
\usepackage{mathtools}
\usepackage{indentfirst}
\usepackage{hyperref}

\usepackage{setspace}
\usepackage{titlesec}
\setlength\parident{1.5em}
\begin{document}
\thispagestyle{empty}

\begin{center}
\ \vspace{-3cm}

\includegraphics[width=0.5\textwidth]{msu.eps}\\
{\scshape Московский государственный университет имени М.~В.~Ломоносова}\\
Факультет вычислительной математики и кибернетики\\
Кафедра системного анализа

\vfill

{\LARGE Отчет по практикуму}

\vspace{1cm}

{\Huge\bfseries <<Нелинейная задача оптимального управления>>}
\end{center}

\vspace{1cm}

\begin{flushright}
  \large
  \textit{Студент 315 группы}\\
  Ю.\,М.~Алексеева

  \vspace{5mm}

  \textit{Руководитель практикума}\\
  м.н.с. И.\,А.~Чистяков
\end{flushright}

\vfill

\begin{center}
Москва, 2025
\end{center}

\newpage
\tableofcontents
\newpage
\section{Постановка задачи}
Рассматривается система из двух обкновенных дифференциальных уравнений:\\

\begin{equation}
\label{eq:1}
\begin{cases}
    \dot{x}_1 = x_2+u_2\\
    \dot{x}_2 = u_1+u_2x_2
\end{cases},
~~t\in[0,T],
\end{equation}
где $x_1, x_2 \in \mathbb{R}, u = (u_1, u_2)'\in\mathbb{R}^2.$ На возможные значения управляющих параметров $u_1, u_2$ наложены следующие ограничения:\\
\begin{itemize}
    \item[1.] либо $0 \leqslant u_1 \leqslant \alpha, ~ u_2 \in [k_1, k_2], ~k_2 \geqslant k_1 > 0,$
    \item[2.] либо $u_1 \in \mathbb{R}, ~u_2\in[k_1, k_2], ~k_2\geqslant k_1 >0.$
\end{itemize}
Задан начальный момент времени $t_0 = 0$ и начальная позиция $x(0): ~x_1(0)=0, x_2(0) = S.$ Необходимо за счет выбора программного управления $u$ и начальной позиции перевести систему из заданной начальной позиции в такую позицию в момент времени $T$, в которой $|x_1(T)-L|\leqslant\varepsilon, ~|x_2(T)|\leqslant\varepsilon.$ На множестве всех реализаций программных управлений, переводящих материальную точку в указанное множество, необходимо решить задачу оптимизации:\\
\begin{equation}
\label{eq:2}
    J = \int_0^Tu^4_1(t)dt \rightarrow \min_{u(\cdot)}.
\end{equation}


\newpage
\section{Теоретические выкладки}
\subsection{Принцип максимума Понтрягина}
Введем функцию Гамильтона-Понтрягина:\\
\begin{equation}
\label{eq:3}
    \overline{H}(\psi, x, u) = \psi_0f_0+ \langle \psi, f\rangle,
\end{equation}
где $f = (f_1(x, u), f_2(x, u)) ~-$ правые части системы $(\ref{eq:1})$.\\
Сформулируем принцип максимума Понтрягина.\\
\newtheorem{теорема}{Теорема}
\newcommand{\const}{\mathrm{const}}
\begin{теорема}
    (ПМП). Пусть $x^*(\cdot), u^*(\cdot)) ~-$ оптимальная пара, тогда существует такая вектор-функция сопряженных переменных $\overline{\psi^*}(\cdot): [t_0, T] \rightarrow \mathbb{R}^3,~ \overline{\psi^*}(t) \neq 0, ~\forall t\in [t_0, T]$ такая, что выполнены следующие условия:
    \begin{itemize}
        \item[1.] Вектор-функция $\overline{\psi}(\cdot)$ удовлетворяет сопряженной системе:
        $$
        \begin{cases}
            \dot{\psi}_0^*(t) = 0,\\
            \left.\dot{\psi^*}(t)= -\dfrac{\partial \overline{H}}{\partial x}\right\vert_{\substack{x=x^*(t) \\ u=u^*(t)}}
        \end{cases}
        $$
        \item[2.] Условие максимума: 
        $$\overline{H}(\psi^*(t), x^*(t), u^*(t)) \stackrel{\text{п.в.} t}{=\joinrel=} \sup_{u\in \mathcal{P}}\overline{H}(x^*(t), \psi^*(t), u) = \overline{M}.$$
        \item[3.] $\psi^*_0 = \const, \overline{M}= \const.$
        \item[4.] Условия трансверсальности:
        $$\psi^*(t_0) \perp T_{x^*(t_0)}\mathcal{X}_0,$$
        $$\psi^*(T) \perp T_{x^*(T)}\mathcal{X}_1,$$
        где $T_{x'}\mathcal{X}' ~-$ касательное подпространство множества $\mathcal{X}'$ в точке $x'.$
    \end{itemize}
\end{теорема}
\subsection{Первичный анализ ПМП}
Запишем функцию Гамильтона-Понтрягина для нашей поставленной задачи, она принимает вид:
$$\overline{H}(\psi, x, u) = \psi_0u_1^4+\psi_1(x_2+u_2)+\psi_2(u_1+u_2x_2).$$
Вычислим $\dot{\psi_1}$ и $\dot{\psi_2}$:
$$\dot{\psi_1} = -\dfrac{\partial \overline{H}}{\partial x_1} = 0,~~ \dot{\psi_2} = -\dfrac{\partial \overline{H}}{\partial x_2} = -\psi_1-\psi_2u_2.$$
Таким образом, сопряженная система имеет вид:
$$
\begin{cases}
\dot{\psi_0} = 0,\\
\dot{\psi_1} = 0,\\
\dot{\psi_2} = -\psi_1-\psi_2u_2,
\end{cases}
\implies
\begin{cases}
\psi_0 = \const,\\
\psi_1 = \const = \psi_1^0,\\
\dot{\psi_2} = -\psi_1-\psi_2u_2.
\end{cases}
$$
Сопряженная система не меняет свой вид в зависимости от значений $\psi_0.$ Поэтому в дальнейшем будем пользоваться ей.
Функция Гамильтона-Понтрягина $\overline{H}$ достигает максимума по переменной $u_1,$ если:
$$\dfrac{\partial\overline{H}}{\partial u_1} = 4\psi_0u_1^3 +\psi_2 = 0 \implies u_1 = -\sqrt[3]{\dfrac{\psi_2}{4\psi_0}}, ~\psi_0 \neq 0.$$
Покажем, что $u_1=-\sqrt{\dfrac{\psi_2}{4\psi_0}},~ \psi_0 \neq0,$ является максимумом функции Гамильтона-Понтрягина. \\
Разложим $\dfrac{\partial\overline{H}}{\partial u_1}$ на множители:
$$4\psi_0u_1^3+\psi_2 = \left(u_1 + \sqrt[3]{\dfrac{\psi_2}{4\psi_0}}\right)\left(4\psi_0u_1^2-u_1\sqrt[3]{16\psi_0^2\psi_2}+\sqrt[3]{4\psi_0\psi_2^2}\right).$$
Оценим вторую скобку:
\begin{multline}
L = \left(4\psi_0 u_1^2 - u_1 \sqrt[3]{16\psi_0^2\psi_2} + \sqrt[3]{4\psi_0\psi_2^2}\right)= \\
= 4\psi_0\left(u_1^2 - u_1 \sqrt[3]{\dfrac{\psi_2}{4\psi_0}} + \sqrt[3]{\dfrac{\psi_2^2}{16\psi_0^2}}\right)= \\
= 4\psi_0\left(\left(u_1 - \dfrac{1}{2}\sqrt[3]{\dfrac{\psi_2}{4\psi_0}}\right)^2 + \dfrac{3}{4}\sqrt[3]{\dfrac{\psi_2^2}{16\psi_0^2}}\right).
\end{multline}
При условии, что $\psi_0 <0$, имеем $L\leqslant0.$ Рассмотрим, что происходит при $L=0 $.
$$L = 0 ~~\Leftrightarrow ~~
\begin{cases}
    u_1 = \dfrac{1}{2}\sqrt[3]{\dfrac{\psi_2}{4\psi_0}},\vspace{5pt}\\
    \dfrac{3}{4}\sqrt[3]{\dfrac{\psi_2^2}{16\psi_0^2}} = 0
\end{cases}
~~\Leftrightarrow ~~
\begin{cases}
    u_1 = 0,\\
    \psi_2 = 0.
\end{cases}
$$
Таким образом, уравнение $4\psi_0u_1^3 +\psi_2 = 0$ имеет три корня, каждый из которых равен 0.\\
Каждый из корней совпадает с $u_1 = -\sqrt[3]{\dfrac{\psi_2}{4\psi_0}}, ~\psi_2=0.$ Значит, кроме $u_1 = -\sqrt[3]{\dfrac{\psi_2}{4\psi_0}}$ других корней нет. Следовательно, осталось сравнить значение производной $\dfrac{\partial\overline{H}}{\partial u_1}$ с нулем и определить поведение функции $\overline{H}$ на интервалах: $\left(-\infty, -\sqrt[3]{\dfrac{\psi_2}{4\psi_0}}\right)$ и $\left(-\sqrt[3]{\dfrac{\psi_2}{4\psi_0}}, +\infty\right)$. \vspace{5pt}\\
Если $u_1\in\left(-\infty, -\sqrt[3]{\dfrac{\psi_2}{4\psi_0}}\right):$
$$u_1+\sqrt[3]{\dfrac{\psi_2}{4\psi_0}} < 0, ~L < 0 \implies \dfrac{\partial\overline{H}}{\partial u_1} > 0 \implies \overline{H} \text{ возрастает.}$$
Если $u_1\in\left(-\sqrt[3]{\dfrac{\psi_2}{4\psi_0}}, +\infty\right):$
$$u_1+\sqrt[3]{\dfrac{\psi_2}{4\psi_0}} > 0, ~L < 0 \implies \dfrac{\partial\overline{H}}{\partial u_1} < 0 \implies \overline{H} \text{ убывает.}$$
Таким образом, $u_1 =  -\sqrt[3]{\dfrac{\psi_2}{4\psi_0}}$, действительно, является максимумом функции Гамильтона-Понтрягина.\\

\newtheorem{определение}{Определение}
\begin{определение}
    Рассматривается линейное неоднородное дифференциальное уравнение:
    $$
    \begin{cases}
        \dot{x}(t) = A(t)x(t) +F(t),\\
        x(t_0) = x^0.
    \end{cases}
    $$
    Пусть $X(t, \tau)~-~$фундаментальная матрица Коши. Тогда решение данного уравнения находится по \textbf{формуле Коши}:
    $$x(t) = X(t, t_0)x^0+\int^t_{t_0}X(t, \tau)F(\tau)d\tau,$$
    где матрица $X(t, \tau)~-~$ решение задачи Коши:
    $$
    \begin{cases}
        \dfrac{\partial X(t, \tau)}{\partial t} = A(t)X(t, \tau),\\
        X(\tau, \tau) = I.
    \end{cases}
    $$
\end{определение}
\textbf{Связь переменных $T, L, S, \varepsilon$}\\
Рассмотрим уравнение:
$$
\begin{cases}
    \dot{x}_2(t) = u_2(t)x_2(t)+u_1(t),\\
    x_2(0) = S.
\end{cases}
$$
Найдем фундаментальную матрицу Коши для рассматриваемого уравнения:
$$
\begin{cases}
    \dfrac{\partial X(t, \tau)}{\partial t} = u_2(t)X(t, \tau),\\
    X(\tau, \tau) = I.
\end{cases}
\implies X(t, \tau) = e^{\int^t_{\tau}u_2(\tau)d\tau}.
$$
Запишем формулу Коши для $x_2(T):$
$$x_2(T) = e^{\int_0^T u_2(t) dt} x_2(0) + e^{\int_0^T u_2(t) dt} \int_0^T u_1(t) e^{-\int_0^t u_2(\tau) d\tau} dt.$$
Учитывая, что $x_2(0)= S,~|x_2(T)|\leqslant\varepsilon, ~0\leqslant u_1\leqslant\alpha, ~k_1\leqslant u_2\leqslant k_2.$
$$k_1\leqslant u_2\leqslant k_2\implies k_1T\leqslant {\int_0^T u_2(t) dt}\leqslant k_2T \implies e^{k_1T}\leqslant 
e^{\int_0^T u_2(t) dt}\leqslant e^{k_2T},$$

\begin{multline*}
k_1 \leqslant u_2 \leqslant k_2 \implies -k_2 t \leqslant -\int_0^t u_2(\tau) d\tau \leqslant -k_1 t 
\implies e^{-k_2 t} \leqslant e^{-\int_0^t u_2(\tau) d\tau} \leqslant e^{-k_1 t} \implies\\
\implies 0 \leqslant u_1 e^{-\int_0^t u_2(\tau) d\tau} \leqslant \alpha e^{-k_1 t}
\implies 0 \leqslant \int_0^T u_1(t) e^{-\int_0^t u_2(\tau) d\tau} dt \leqslant -\dfrac{\alpha}{k_1} \left( e^{-k_1 T} - 1 \right) \implies\\
\implies 0 \leqslant e^{\int_0^T u_2(t) dt} \int_0^T u_1(t) e^{-\int_0^t u_2(\tau) d\tau} dt \leqslant \dfrac{\alpha}{k_1} \left( e^{k_2 T} - e^{(k_2 - k_1) T} \right).
\end{multline*}

Значит,

\begin{multline*}
Se^{k_1 T} \leqslant e^{\int_0^T u_2(t) dt} x_2(0) + e^{\int_0^T u_2(t) dt} \int_0^T u_1(t) e^{-\int_0^t u_2(\tau) d\tau} dt 
\leqslant \dfrac{\alpha}{k_1} \left( e^{k_2 T} - e^{(k_2 - k_1) T} \right) + Se^{k_2 T} \implies\\
\implies  -\varepsilon \leqslant Se^{k_1 T} \leqslant x_2(T)  \leqslant \dfrac{\alpha}{k_1} \left( e^{k_2 T} - e^{(k_2 - k_1) T} \right) + Se^{k_2 T}\leqslant \varepsilon.
\end{multline*}
Таким образом, получаем итоговое неравенство на связь интересующих нас переменных:
$$-\varepsilon \leqslant Se^{k_1 T} \leqslant \dfrac{\alpha}{k_1} \left( e^{k_2 T} - e^{(k_2 - k_1) T} \right) + Se^{k_2 T}\leqslant \varepsilon$$








\subsubsection{Задача 1}
$\label{subsub:1}$
Напомним, что $0 \leqslant u_1 \leqslant \alpha, ~ u_2 \in [k_1, k_2], ~k_2 \geqslant k_1 > 0.$\\
\textbf{Нормальный случай}\\
Здесь $\psi_0\neq0.$ Не ограничивая общности, можно считать $\psi_0 = -\dfrac{1}{4}.$\\
Для максимизации функции Гамильтона-Понтрягина выбираем следующие управления:
$$
u_1^* = 
\begin{cases}
    \alpha, & \psi_2 > \alpha^3, \\
    \sqrt[3]{\psi_2}, & 0 \leqslant \psi_2 \leqslant \alpha^3, \\
    0, & \psi_2 < 0,
\end{cases}
~~~~
u_2^* = 
\begin{cases}
    k, & \psi_1 + \psi_2 x_2 > 0, \\
    [-k, k], & \psi_1 + \psi_2 x_2 = 0, \\
    -k, & \psi_1 + \psi_2 x_2 < 0.
\end{cases}
$$
\textbf{Анормальный случай}\\
Здесь $\psi_0=0.$\\
Для максимизации функции Гамильтона-Понтрягина выбираем следующие управления:
$$
u_1^* = 
\begin{cases}
    \alpha, &\psi_2 >0, \\
    [0, \alpha], &\psi_2=0,\\
    0, &\psi_2 < 0,
\end{cases}
~~~~
u_2^* = 
\begin{cases}
    k, &\psi_1 +\psi_2x_2>0, \\
    [-k, k], &\psi_1 +\psi_2x_2=0,\\
    -k, &\psi_1 +\psi_2x_2<0.
\end{cases}
$$
\textbf{Условия трансверсальности}\\
Рассмотрим условия трансверсальности для поставленной задачи. \\
Исследуем достижимость точки $(L, 0)$. \\
Если $u_1\geqslant 0,~ x_2(T) = 0,$ то запишем формулу Коши для $x_2(T)$:
$$x_2(T) = e^{\int_0^T u_2(t) dt} x_2(0) + e^{\int_0^T u_2(t) dt} \int_0^T u_1(t) e^{-\int_0^t u_2(\tau) d\tau} dt = 0 \implies S + \int_0^T u_1(t) e^{-\int_0^t u_2(\tau) d\tau} dt = 0.$$
Таким образом, рассмотрим 3 случая: $S> 0,S = 0,  S <0.$\\
\begin{itemize}
    \item [1.] $S > 0$:
    $$u_1(t) \geqslant 0, e^{-\int_0^t u_2(\tau) d\tau} > 0 \implies x_2(T) > 0 \implies \text{противоречие, точка недостижима.}$$
    $$u_1(t) \in \mathbb{R}, e^{-\int_0^t u_2(\tau) d\tau} > 0 \implies \text{точка достижима.} $$
    Управление в этом случае должно удовлетворять уравнению:
    $$S + \int_0^T u_1(t) e^{-\int_0^t u_2(\tau) d\tau} dt = 0.$$
    \item[2.] $S=0:$
    $$u_1(t) \geqslant 0, e^{-\int_0^t u_2(\tau) d\tau} > 0 \implies u_1(t) =  0, ~ \forall t\in[0, T].$$
    Значит, $$\dot{x}_1 = 0 + 0, ~x(0) = 0 \implies x_1(t) = 0, ~\forall t\in[0,T].$$
    Пришли к противоречию. В этом случае, точка недостижима.
    \item [3.] $S<0:$
    Точка достижима и управление в этом случае должно удовлетворять уравнению:
    $$S + \int_0^T u_1(t) e^{-\int_0^t u_2(\tau) d\tau} dt = 0.$$
\end{itemize}
Чтобы обобщить все эти случаи, для простоты в качестве конечного множества можно рассматривать любую окрестность точки $(L, 0)$. В качестве примера рассмотрим случай, когда конечной точкой множества является квадрат.\\
Множество $\mathcal{X}_0$ является одноточечным множеством, поэтому условие трансверсальности на левом конце будет выполнено для любой функции $\psi(t).$\\
Множество $\mathcal{X}_1$ представляет собой квадрат $[L-\varepsilon, L+\varepsilon]\times [-\varepsilon, \varepsilon].$ Таким образом, условие трансверсальности для правого конца (при нормировке вектора $\psi(t)$) можно записать в виде:
$$
\psi(t) = 
\begin{cases}
    (1, 0), &\text{если}  ~~x_1(T) = \varepsilon, ~x_2(T)\in[L-\varepsilon, L+\varepsilon
    ],\\
    (-1, 0), &\text{если}  ~~x_1(T) = -\varepsilon, ~x_2(T)\in[L-\varepsilon, L+\varepsilon
    ],\\
    (0, 1), &\text{если}  ~~x_2(T) = L + \varepsilon, ~x_1(T)\in[-\varepsilon, \varepsilon
    ],\\
    (0, -1), &\text{если}  ~~x_2(T) = L- \varepsilon, ~x_1(T)\in[-\varepsilon, \varepsilon
    ].\\
\end{cases}
$$
\subsubsection{Задача 2}
Напомним, что $u_1 \in \mathbb{R}, ~u_2\in[k_1, k_2], ~k_2\geqslant k_1 >0.$\\
\textbf{Нормальный случай}
Здесь $\psi_0\neq0.$ Не ограничивая общности, можно считать $\psi_0 = -\dfrac{1}{4}.$\\
Для максимизации функции Гамильтона-Понтрягина выбираем следующие управления:
$$
u_1^* = 
\sqrt[3]{\psi_2},
~~~~
u_2^* = 
\begin{cases}
    k, & \psi_1 + \psi_2 x_2 > 0, \\
    [-k, k], & \psi_1 + \psi_2 x_2 = 0, \\
    -k, & \psi_1 + \psi_2 x_2 < 0.
\end{cases}
$$
\textbf{Анормальный случай}\\
Анормальный случай невозможен, так как при $\psi_0= 0$ управление не ограниченно $(u_1(t) = \infty).$\\
\textbf{Условия трансверсальности}\\
Условия трансверсальности такие же, как в пункте $\ref{subsub:1}.$
\newpage
\begin{thebibliography}{99}
\bibitem{1} Чистяков И.А., Паршиков М.В. Лекции по Оптимальному Управлению
\bibitem{2} Арутюнов А.В. Лекции по Многозначному Анализу.
\end{thebibliography}
\end{document}
